%
% File main.tex
%
% NLP course reproducibility-checklist template, Warsaw University of Technology.
% Prepared by Anna Wróblewska, adapted from the ACL 2015 conference
% formatting template (acl2015.sty), originally by the ACL 2015
% organizers and earlier ACL/EACL format authors (see acl2015.pdf
% for the full credits chain).

\documentclass[11pt]{article}
\usepackage{acl2015}
\usepackage{times}
\usepackage{url}
\usepackage{latexsym}
\usepackage{xcolor}

%\setlength\titlebox{5cm}

% You can expand the title box if you need extra space
% to show all the authors. Please do not make the title box
% smaller than 5cm (the original size); we will check this
% in the camera-ready version and ask you to change it back.


\title{Reproducibility Appendix\\Project Report for NLP Course, Winter 2026/27}
% https://www.overleaf.com/project/63440c16442dba647b52a1ae
\author{Author 1 (submitted by) \\
  Album number \\
  {\tt email@pw.edu.pl} \\
  \And
  Author 2 \\
  Album number \\
  {\tt email@pw.edu.pl} \\
  \AND
  supervisor: Anna Wróblewska\\
  Warsaw University of Technology \\
  {\tt anna.wroblewska1@pw.edu.pl}
  }

\date{}
\begin{document}
\maketitle

{\small\em Template: Anna Wr\'oblewska (Warsaw University of Technology), adapted from the ACL 2015 conference formatting template.}
\vspace{4pt}


\section*{Reproducibility checklist}

% This list is based on the paper: Ian Magnusson, Noah A. Smith, and Jesse Dodge. 2023. Reproducibility in NLP: What Have We Learned from the Checklist?. In Findings of the Association for Computational Linguistics: ACL 2023, pages 12789–12811, Toronto, Canada. Association for Computational Linguistics.

Overall results:
\begin{itemize}
    \item MODEL DESCRIPTION -- A clear description of the mathematical setting, algorithm, and/or model
    \item LINK TO CODE -- A link to a downloadable source code, with specification of all dependencies, including external libraries
    \item INFRASTRUCTURE -- A description of the computing infrastructure used
    \item RUNTIME PARAMETERS -- Average runtime for each approach
    \item PARAMETERS -- The number of parameters in each model
    \item VALIDATION PERFORMANCE -- Corresponding validation performance for each reported test result
    \item METRICS -- Explanation of evaluation metrics used, with links to code
    \item \textbf{SEEDS} -- Random seeds used for each reported run, so the exact run can be reproduced
    \item \textbf{EXACT COMMANDS} -- The exact command(s) or script used to produce the main reported result, with library/package versions
\end{itemize}

Multiple Experiments:
\begin{itemize}
    \item NO TRAINING EVAL RUNS -- The exact number of training and evaluation runs
    \item HYPER BOUND -- Bounds for each hyperparameter
    \item HYPER BEST CONFIG -- Hyperparameter configurations for best-performing models
    \item HYPER SEARCH -- Number of hyperparameter search trials
    \item HYPER METHOD -- The method of choosing hyperparameter values (e.g., uniform sampling, manual tuning, etc.) and the criterion used to select among them (e.g., accuracy)
    \item EXPECTED PERF -- Summary statistics of the results (e.g., mean, variance, error bars, etc.)
\end{itemize}

Datasets -- utilized in the experiments and/or the created ones:
\begin{itemize}
    \item DATA STATS -- Relevant statistics, such as the number of examples
    \item DATA SPLIT -- Details of train/validation/test splits
    \item DATA PROCESSING -- Explanation of any data that were excluded and all preprocessing steps
    \item DATA DOWNLOAD -- A link to a downloadable version of the data
    \item NEW DATA DESCRIPTION --  For new data collected, a complete description of the data collection process, such as instructions to annotators and methods for quality control
    \item DATA LANGUAGES -- For natural language data, the name of the language(s)
\end{itemize}

\end{document}
